\documentclass[format=sigplan, review=false, screen=true]{acmart}
\usepackage[T1]{fontenc}
\usepackage[utf8]{inputenc}
\usepackage{fontspec}
\setmonofont[
  Contextuals={Alternate},
  Scale=0.9
]{Fira Code}
\makeatletter
\def\verbatim@nolig@list{}
\makeatother

\usepackage{booktabs} % For formal tables
\usepackage{bussproofs}
\usepackage{graphics}
\usepackage{xcolor}

\usepackage{minted}
\usemintedstyle{friendly}
\setminted{
    frame=lines,
    framesep=2mm,
    baselinestretch=1.2,
    %bgcolor=lightgray,
    fontsize=\footnotesize
    }

\usepackage{filecontents}

\usepackage[ruled]{algorithm2e} % For algorithms
\renewcommand{\algorithmcfname}{ALGORITHM}
\SetAlFnt{\small}
\SetAlCapFnt{\small}
\SetAlCapNameFnt{\small}
\SetAlCapHSkip{0pt}
\IncMargin{-\parindent}

\usepackage{listings}
\lstdefinestyle{mystyle}{
    %backgroundcolor=\color{backcolour},   
    commentstyle=\color{gray},
    keywordstyle=\color{blue},
    numberstyle=\tiny,
    stringstyle=\color{purple},
    basicstyle=\small\tt,
    breakatwhitespace=false,         
    breaklines=true,                 
    captionpos=b,                    
    keepspaces=true,                 
    %numbers=left,                    
    numbersep=5pt,                  
    showspaces=false,                
    showstringspaces=false,
    showtabs=false,
    frame=single,
    xleftmargin=1em,
    xrightmargin=1em,
    frame=shadowbox,
    rulesepcolor=\color{gray},
    tabsize=2
}
 
\lstset{style=mystyle,
  literate=
  {->} {$\to$} 2
  {<-} {$\leftarrow$} 2
  {=>} {$\Rightarrow$} 2
  {forall} {$\forall$} 1
  {exists} {$\exists$} 1
  {phi} {$\varphi$} 1
  {rho} {$\rho$} 1
  {kappa} {$\kappa$} 1
  {$nu$} {$\nu$} 1
  {$mu$} {$\mu$} 1
  {gamma} {$\gamma$} 1
  {subsetX} {$\subset$} 1
  {~>} {$\rightsquigarrow$} 2
  {<~>} {$\leftrightsquigarrow$} 3
  {elem} {$\in$} 1
}

% Metadata Information
\acmJournal{TWEB}
\acmVolume{9}
\acmNumber{4}
\acmArticle{39}
\acmYear{2018}
\acmMonth{3}
\copyrightyear{2018}
%\acmArticleSeq{9}

% Copyright
%\setcopyright{acmcopyright}
\setcopyright{acmlicensed}
%\setcopyright{rightsretained}
%\setcopyright{usgov}
%\setcopyright{usgovmixed}
%\setcopyright{cagov}
%\setcopyright{cagovmixed}

% DOI
\acmDOI{0000001.0000001}

% Paper history
\received{June 2018}
%\received[revised]{March 2009}
%\received[accepted]{June 2009}

% Document starts
\begin{document}
% Title portion. Note the short title for running heads
\title[Ghosts of Departed Proofs]{Functional Pearl: Ghosts of Departed Proofs}

\author{Matt Noonan}
\orcid{1234-5678-9012-3456}
\affiliation{%
  \institution{Kataskeue LLC}
%  \streetaddress{Esty St}
  \city{Ithaca}
  \state{NY}
  \postcode{14850}
  \country{USA}}
\email{mnoonan@kataskeue.com}


\begin{abstract}

  We present a simple technique that allows library authors to
  control how APIs are used.

\end{abstract}


%
% The code below should be generated by the tool at
% http://dl.acm.org/ccs.cfm
% Please copy and paste the code instead of the example below.
%
\begin{CCSXML}
<ccs2012>
 <concept>
  <concept_id>10010520.10010553.10010562</concept_id>
  <concept_desc>Computer systems organization~Embedded systems</concept_desc>
  <concept_significance>500</concept_significance>
 </concept>
 <concept>
  <concept_id>10010520.10010575.10010755</concept_id>
  <concept_desc>Computer systems organization~Redundancy</concept_desc>
  <concept_significance>300</concept_significance>
 </concept>
 <concept>
  <concept_id>10010520.10010553.10010554</concept_id>
  <concept_desc>Computer systems organization~Robotics</concept_desc>
  <concept_significance>100</concept_significance>
 </concept>
 <concept>
  <concept_id>10003033.10003083.10003095</concept_id>
  <concept_desc>Networks~Network reliability</concept_desc>
  <concept_significance>100</concept_significance>
 </concept>
</ccs2012>
\end{CCSXML}

\ccsdesc[500]{Computer systems organization~Embedded systems}
\ccsdesc[300]{Computer systems organization~Redundancy}
\ccsdesc{Computer systems organization~Robotics}
\ccsdesc[100]{Networks~Network reliability}

%
% End generated code
%


\keywords{Wireless sensor networks, media access control,
multi-channel, radio interference, time synchronization}




\maketitle

% The default list of authors is too long for headers.
\renewcommand{\shortauthors}{M. Noonan}

%\section{Introduction}


As a new technology, Wireless Sensor Networks (WSNs) has a wide
range of applications \cite{Culler-01, Bahl-02, Akyildiz-01}, including
environment monitoring, smart buildings, medical care, industrial and
military applications. Among them, a recent trend is to develop
commercial sensor networks that require pervasive sensing of both
environment and human beings, for example, assisted living
\cite{Akyildiz-02, Harvard-01,CROSSBOW} and smart homes
\cite{Harvard-01, Adya-01,CROSSBOW}.
% quote
\begin{quote}
  ``For these applications, sensor devices are incorporated into human
  cloths \cite{Natarajan-01, Zhou-06, Bahl-02, Adya-01} for monitoring
  health related information like EKG readings, fall detection, and
  voice recognition''.
\end{quote}
While collecting all these multimedia information
\cite{Akyildiz-02} requires a high network throughput, off-the-shelf
sensor devices only provide very limited bandwidth in a single
channel: 19.2\,Kbps in MICA2 \cite{Bahl-02} and 250\,Kbps in MICAz.

In this article, we propose MMSN, abbreviation for Multifrequency
Media access control for wireless Sensor Networks. The main
contributions of this work can be summarized as follows.
% itemize
\begin{itemize}
\item To the best of our knowledge, the MMSN protocol is the first
multifrequency MAC protocol especially designed for WSNs, in which
each device is equipped with a single radio transceiver and
the MAC layer packet size is very small.
\item Instead of using pairwise RTS/CTS frequency negotiation
\cite{Adya-01, Culler-01, Tzamaloukas-01, Zhou-06},
we propose lightweight frequency assignments, which are good choices
for many deployed comparatively static WSNs.
\item We develop new toggle transmission and snooping techniques to
enable a single radio transceiver in a sensor device to achieve
scalable performance, avoiding the nonscalable ``one
control channel + multiple data channels'' design \cite{Natarajan-01}.
\end{itemize}

% Head 1
\section{MMSN Protocol}

% Head 2
\subsection{Frequency Assignment}

We propose a suboptimal distribution to be used by each node, which is
easy to compute and does not depend on the number of competing
nodes. A natural candidate is an increasing geometric sequence, in
which
% Numbered Equation
\begin{equation}
\label{eqn:01}
P(t)=\frac{b^{\frac{t+1}{T+1}}-b^{\frac{t}{T+1}}}{b-1},
\end{equation}
where $t=0,{\ldots}\,,T$, and $b$ is a number greater than $1$.

In our algorithm, we use the suboptimal approach for simplicity and
generality. We need to make the distribution of the selected back-off
time slice at each node conform to what is shown in
Equation~\eqref{eqn:01}. It is implemented as follows: First, a random
variable $\alpha$ with a uniform distribution within the interval $(0,
1)$ is generated on each node, then time slice $i$ is selected
according to the following equation:
% Unnumbered Equation
\[
i=\lfloor(T+1)\log_b[\alpha(b-1)+1]\rfloor.
\]
It can be easily proven that the distribution of $i$ conforms to Equation
(\ref{eqn:01}).

So protocols \cite{Bahl-02, Culler-01,Zhou-06,Adya-01,
Tzamaloukas-01, Akyildiz-01} that use RTS/CTS
controls\footnote{RTS/CTS controls are required to be implemented by
802.11-compliant devices. They can be used as an optional mechanism
to avoid Hidden Terminal Problems in the 802.11 standard and
protocols based on those similar to \cite{Akyildiz-01} and
\cite{Adya-01}.} for frequency negotiation and reservation are not
suitable for WSN applications, even though they exhibit good
performance in general wireless ad-hoc
networks.

% Head 3
\subsubsection{Exclusive Frequency Assignment}


In exclusive frequency assignment, nodes first exchange their IDs
among two communication hops so that each node knows its two-hop
neighbors' IDs. In the second broadcast, each node beacons all
neighbors' IDs it has collected during the first broadcast period.

% Head 4
\paragraph{Eavesdropping}

Even though the even selection scheme leads to even sharing of
available frequencies among any two-hop neighborhood, it involves a
number of two-hop broadcasts. To reduce the communication cost, we
propose a lightweight eavesdropping scheme.

\subsection{Basic Notations}

As Algorithm~\ref{alg:one} states, for each frequency
number, each node calculates a random number (${\textit{Rnd}}_{\alpha}$) for
itself and a random number (${\textit{Rnd}}_{\beta}$) for each of its two-hop
neighbors with the same pseudorandom number generator.

% Algorithm
\begin{algorithm}[t]
\SetAlgoNoLine
\KwIn{Node $\alpha$'s ID ($ID_{\alpha}$), and node $\alpha$'s
neighbors' IDs within two communication hops.}
\KwOut{The frequency number ($FreNum_{\alpha}$) node $\alpha$ gets assigned.}
$index$ = 0; $FreNum_{\alpha}$ = -1\;
\Repeat{$FreNum_{\alpha} > -1$}{
        $Rnd_{\alpha}$ = Random($ID_{\alpha}$, $index$)\;
        $Found$ = $TRUE$\;
        \For{each node $\beta$ in $\alpha$'s two communication hops
    }{
      $Rnd_{\beta}$ = Random($ID_{\beta}$, $index$)\;
      \If{($Rnd_{\alpha} < Rnd_{\beta}$) \text{or} ($Rnd_{\alpha}$ ==
          $Rnd_{\beta}$ \text{and} $ID_{\alpha} < ID_{\beta}$)\;
      }{
        $Found$ = $FALSE$; break\;
      }
        }
     \eIf{$Found$}{
           $FreNum_{\alpha}$ = $index$\;
         }{
           $index$ ++\;
     }
      }
\caption{Frequency Number Computation}
\label{alg:one}
\end{algorithm}


Bus masters are divided into two disjoint sets, $\mathcal{M}_{RT}$
and $\mathcal{M}_{NRT}$.
% description
\begin{description}
\item[RT Masters]
$\mathcal{M}_{RT}=\{ \vec{m}_{1},\dots,\vec{m}_{n}\}$ denotes the
$n$ RT masters issuing real-time constrained requests. To model the
current request issued by an $\vec{m}_{i}$ in $\mathcal{M}_{RT}$,
three parameters---the recurrence time $(r_i)$, the service cycle
$(c_i)$, and the relative deadline $(d_i)$---are used, with their
relationships.
\item[NRT Masters]
$\mathcal{M}_{NRT}=\{ \vec{m}_{n+1},\dots,\vec{m}_{n+m}\}$ is a set
of $m$ masters issuing nonreal-time constrained requests. In our
model, each $\vec{m}_{j}$ in $\mathcal{M}_{NRT}$ needs only one
parameter, the service cycle, to model the current request it
issues.
\end{description}

Here, a question may arise, since each node has a global ID. Why
don't we just map nodes' IDs within two hops into a group of
frequency numbers and assign those numbers to all nodes within two
hops?

\section{Simulator}
\label{sec:sim}

If the model checker requests successors of a state which are not
created yet, the state space uses the simulator to create the
successors on-the-fly. To create successor states the simulator
conducts the following steps.
% enumerate
\begin{enumerate}
\item Load state into microcontroller model.
\item Determine assignments needed for resolving nondeterminism.
\item For each assignment.
      \begin{enumerate}
      \item either call interrupt handler or simulate effect of next instruction, or
      \item evaluate truth values of atomic propositions.
      \end{enumerate}
\item Return resulting states.
\end{enumerate}
Figure~\ref{fig:one} shows a typical microcontroller C program that
controls an automotive power window lift. The program is one of the
programs used in the case study described in Section~\ref{sec:sim}.
At first sight, the programs looks like an ANSI~C program. It
contains function calls, assignments, if clauses, and while loops.
% Figure
\begin{figure}
  \includegraphics{mouse}
  \caption{Code before preprocessing.}
  \label{fig:one}
\end{figure}

\subsection{Problem Formulation}

The objective of variable coalescence-based offset assignment is to find
both the coalescence scheme and the MWPC on the coalesced graph. We start
with a few definitions and lemmas for variable coalescence.

% Enunciations
\begin{definition}[Coalesced Node (C-Node)]A C-node is a set of
live ranges (webs) in the AG or IG that are coalesced. Nodes within the same
C-node cannot interfere with each other on the IG. Before any coalescing is
done, each live range is a C-node by itself.
\end{definition}

\begin{definition}[C-AG (Coalesced Access Graph)]The C-AG is the access
graph after node coalescence, which is composed of all C-nodes and C-edges.
\end{definition}

\begin{lemma}
The C-MWPC problem is NP-complete.
\end{lemma}
\begin{proof} C-MWPC can be easily reduced to the MWPC problem assuming a
coalescence graph without any edge or a fully connected interference graph.
Therefore, each C-node is an uncoalesced live range after value separation
and C-PC is equivalent to PC. A fully connected interference graph is made
possible when all live ranges interfere with each other. Thus, the C-MWPC
problem is NP-complete.
\end{proof}

\begin{lemma}[Lemma Subhead]The solution to the C-MWPC problem is no
worse than the solution to the MWPC.
\end{lemma}
\begin{proof}
Simply, any solution to the MWPC is also a solution to the
C-MWPC. But some solutions to C-MWPC may not apply to the MWPC (if any
coalescing were made).
\end{proof}

\section{Performance Evaluation}

During all the experiments, the Geographic Forwarding (GF) by Akyildiz
et al.~\shortcite{Akyildiz-01} routing protocol is used. GF exploits
geographic information of nodes and conducts local data-forwarding to
achieve end-to-end routing. Our simulation is configured according to
the settings in Table~\ref{tab:one}. Each run lasts for 2 minutes and
repeated 100 times. For each data value we present in the results, we
also give its 90\% confidence interval.

% Table
\begin{table}%
\caption{Simulation Configuration}
\label{tab:one}
\begin{minipage}{\columnwidth}
\begin{center}
\begin{tabular}{ll}
  \toprule
  TERRAIN\footnote{This is a table footnote. This is a
    table footnote. This is a table footnote.}   & (200m$\times$200m) Square\\
  Node Number     & 289\\
  Node Placement  & Uniform\\
  Application     & Many-to-Many/Gossip CBR Streams\\
  Payload Size    & 32 bytes\\
  Routing Layer   & GF\\
  MAC Layer       & CSMA/MMSN\\
  Radio Layer     & RADIO-ACCNOISE\\
  Radio Bandwidth & 250Kbps\\
  Radio Range     & 20m--45m\\
  \bottomrule
\end{tabular}
\end{center}
\bigskip\centering
\footnotesize\emph{Source:} This is a table
 sourcenote. This is a table sourcenote. This is a table
 sourcenote.

 \emph{Note:} This is a table footnote.
\end{minipage}
\end{table}%


\section{Conclusions}

In this article, we develop the first multifrequency MAC protocol for
WSN applications in which each device adopts a
single radio transceiver. The different MAC design requirements for
WSNs and general wireless ad-hoc networks are
compared, and a complete WSN multifrequency MAC design (MMSN) is
put forth. During the MMSN design, we analyze and evaluate different
choices for frequency assignments and also discuss the nonuniform
back-off algorithms for the slotted media access design.

% Start of "Sample References" section

\section{Typical References in New ACM Reference Format}
A paginated journal article \cite{Abril07}, an enumerated
journal article \cite{Cohen07}, a reference to an entire issue \cite{JCohen96},
a monograph (whole book) \cite{Kosiur01}, a monograph/whole book in a series (see 2a in spec. document)
\cite{Harel79}, a divisible-book such as an anthology or compilation \cite{Editor00}
followed by the same example, however we only output the series if the volume number is given
\cite{Editor00a} (so Editor00a's series should NOT be present since it has no vol. no.),
a chapter in a divisible book \cite{Spector90}, a chapter in a divisible book
in a series \cite{Douglass98}, a multi-volume work as book \cite{Knuth97},
an article in a proceedings (of a conference, symposium, workshop for example)
(paginated proceedings article) \cite{Andler79}, a proceedings article
with all possible elements \cite{Smith10}, an example of an enumerated
proceedings article \cite{VanGundy07},
an informally published work \cite{Harel78}, a doctoral dissertation \cite{Clarkson85},
a master's thesis: \cite{anisi03}, an online document / world wide web
resource \cite{Thornburg01, Ablamowicz07, Poker06}, a video game (Case 1) \cite{Obama08} and (Case 2) \cite{Novak03}
and \cite{Lee05} and (Case 3) a patent \cite{JoeScientist001},
work accepted for publication \cite{rous08}, 'YYYYb'-test for prolific author
\cite{SaeediMEJ10} and \cite{SaeediJETC10}. Other cites might contain
'duplicate' DOI and URLs (some SIAM articles) \cite{Kirschmer:2010:AEI:1958016.1958018}.
Boris / Barbara Beeton: multi-volume works as books
\cite{MR781536} and \cite{MR781537}.

A couple of citations with DOIs: \cite{2004:ITE:1009386.1010128,
  Kirschmer:2010:AEI:1958016.1958018}.

Online citations: \cite{TUGInstmem, Thornburg01, CTANacmart}.

% Appendix
\appendix
\section{Switching Times}

In this appendix, we measure the channel switching time of Micaz
\cite{CROSSBOW} sensor devices.  In our experiments, one mote
alternatingly switches between Channels~11 and~12. Every time after
the node switches to a channel, it sends out a packet immediately and
then changes to a new channel as soon as the transmission is finished.
We measure the number of packets the test mote can send in 10 seconds,
denoted as $N_{1}$. In contrast, we also measure the same value of the
test mote without switching channels, denoted as $N_{2}$. We calculate
the channel-switching time $s$ as
\begin{displaymath}%
s=\frac{10}{N_{1}}-\frac{10}{N_{2}}.
\end{displaymath}%
By repeating the experiments 100 times, we get the average
channel-switching time of Micaz motes: 24.3\,$\mu$s.

\section{Supplementary Materials}


\begin{printonly}
  See the supplementary materials in the online version
\end{printonly}

\begin{screenonly}
\subsection{This is an Example of Appendix Subsection Head}

Channel-switching time is measured as the time length it takes for
motes to successfully switch from one channel to another. This
parameter impacts the maximum network throughput, because motes
cannot receive or send any packet during this period of time, and it
also affects the efficiency of toggle snooping in MMSN, where motes
need to sense through channels rapidly.

By repeating experiments 100 times, we get the average
channel-switching time of Micaz motes: 24.3 $\mu$s. We then conduct
the same experiments with different Micaz motes, as well as
experiments with the transmitter switching from Channel 11 to other
channels. In both scenarios, the channel-switching time does not have
obvious changes. (In our experiments, all values are in the range of
23.6 $\mu$s to 24.9 $\mu$s.)

\subsection{Appendix Subsection Head}

The primary consumer of energy in WSNs is idle listening. The key to
reduce idle listening is executing low duty-cycle on nodes. Two
primary approaches are considered in controlling duty-cycles in the
MAC layer.

\end{screenonly}

\begin{acks}

The authors would like to thank Dr. Maura Turolla of Telecom
Italia for providing specifications about the application scenario.

The work is supported by the \grantsponsor{GS501100001809}{National
  Natural Science Foundation of
  China}{http://dx.doi.org/10.13039/501100001809} under Grant
No.:~\grantnum{GS501100001809}{61273304\_a}
and~\grantnum[http://www.nnsf.cn/youngscientists]{GS501100001809}{Young
  Scientists' Support Program}.


\end{acks}

% Bibliography
\bibliographystyle{ACM-Reference-Format}
\bibliography{sample-bibliography}


\section{Introduction}

\subsection{Encoding with Universals}

It is a theorem of both classical and constructive logics that
\[\forall t.~ (\forall s. \varphi(s) \Rightarrow t) \Rightarrow t \equiv \exists c.~ \varphi(c)\]

\section{Warmup: Not quite dependent types}

\begin{filecontents*}{ex1.hs}  
{-# LANGUAGE RankNTypes #-}
module Sized
  (Size, the, sZipWith, sizing, align) where

newtype Size n a = Size a

the :: Size n a -> a
the (Size x) = x

sZipWith :: (a -> b -> c)
         -> Size n [a]
         -> Size n [b]
         -> Size n [c]
sZipWith f xs ys =
  Size (zipWith f (the xs) (the ys))
  
sizing :: [a] -> (forall n. Size n [a] -> t) -> t
sizing xs k = k (Size xs)

align :: Size n [a] -> [b] -> Maybe (Size n [b])
align xs ys = if length (the xs) == length ys
              then Just (Size ys)
              else Nothing

\end{filecontents*}

\begin{filecontents*}{ex2.hs}
import Sized

dot :: Size n [Double] -> Size n [Double] -> Double
dot xs ys = sum (the $ sZipWith (*) xs ys)

main :: IO ()
main = do
  xs <- readLn
  ys <- readLn
  sizing xs $ \xs' -> do
    case align xs' ys of
      Nothing  -> putStrLn "Size mismatch!"
      Just ys' -> print (dot xs' ys')
\end{filecontents*}

\begin{figure}
    \inputminted{haskell}{ex1.hs}
  \caption{A small module defining a type for lists with a known length.}
\end{figure}
\begin{figure}
    \inputminted{haskell}{ex2.hs}
    \caption{A user-defined dot product function that can only be used on same-sized lists,
    and a usage example.}
\end{figure}


\begin{lstlisting}[language=Haskell]
  norm2 :: [Double] -> Double
  norm2 xs = sizing xs (\v -> v `dot` v)
\end{lstlisting}

\begin{lstlisting}[language=Haskell]
  sizing xs $ \xs' -> 
    case align xs' ys of
      Just ys' -> (xs' `dot` ys') / (xs' `dot` xs')
      Nothing  -> 17
\end{lstlisting}

Despite  • appearences, the phantom type parameter $n$ does not really represent the vector's length
{\em per se}. Instead, we propose to think of \texttt{Size n} as a predicate, and values of
type \texttt{Size n [a]} should be thought of as ``lists of type \texttt{[a]}, equipped with a proof
that they satisfy \texttt{Size n}''. Critically, this proof has no run-time impact: it is trapped in
the phantom type parameter.

This approach gives us a straightforward way to interpret the type signatures from example ***:

\begin{lstlisting}[language=Haskell]
  -- You can take the dot product of two lists, if you have proven
  -- that they have the same Size n.
  dot :: Size n [Double] -> Size n [Double] -> Double

  -- When you map a function over a list of Size n, the
  -- result will also have Size n.
  smap :: (a -> b) -> Size n [a] -> Size n [b]

  -- For any list, there is some n such that Size n is true.
  sizing :: [a] -> (forall n. Size n [a] -> t) -> t

  -- Given a list of Size n, we may be able to prove that
  -- another list also has Size n.
  align :: Size n [a] -> [b] -> Maybe (Size n [b])
\end{lstlisting}

As we attach increasingly sophisticated information into the phantom types, it becomes useful to
have a uniform method for \emph{forgetting} all of the ornamentation, revealing the normal
value underneath.

\begin{filecontents*}{theTC.hs}
class The d a | d -> a where
    the :: d -> a
    default the :: Coercible d a => d -> a
    the = coerce
    
instance The (Size n a) a
\end{filecontents*}

\begin{figure}
  \inputminted{haskell}{theTC.hs}
  \caption{The \texttt{The} typeclass, for dropping ghosts
    from a type. The default instance should always be used,
    so new instances can be created with an empty
    \texttt{instance} declaration.}
\end{figure}

\section{Case Study \#1: Sorted lists}

\begin{filecontents*}{named.hs}
module Named (Named, name) where

import The

newtype Named name a = Named a
instance The (Named name a) a

name :: a -> (forall name. Named name a -> t) -> t
name x k = k (coerce x)
\end{filecontents*}

\begin{filecontents*}{ordered.hs}
module Sorted
  (Named, SortedBy, sortBy, mergeBy) where   

import The
import Named

import qualified Data.List as L
import qualified Data.List.Utils as U

newtype SortedBy o a = SortedBy a
instance The (SortedBy o a) a
  
sortBy :: Named comp (a -> a -> Ordering)
       -> [a]
       -> SortedBy comp [a]
sortBy comp xs = coerce (L.sortBy (the comp) xs)

mergeBy :: Named comp (a -> a -> Ordering)
        -> SortedBy comp [a]
        -> SortedBy comp [a]
        -> SortedBy comp [a]
mergeBy comp xs ys =
  coerce (U.mergeBy (the comp) (the xs) (the ys))        
\end{filecontents*}

\begin{filecontents*}{usageO.hs}
import Sorted
import Named
main = do
  xs <- readLn :: IO Int
  ys <- readLn
  name (>) $ \gt -> do
    let xs' = sortBy gt xs
        ys' = sortBy gt ys
    print (the xs', the ys', the (mergeBy gt xs' ys'))
\end{filecontents*}

\begin{figure}
  \inputminted{haskell}{ordered.hs}
  \caption{A module for working with lists that have been sorted by an arbitrary
    comparator.}
\end{figure}


\begin{figure}
  \inputminted{haskell}{named.hs}
  \caption{A module for attaching ghostly names to values.}
\end{figure}

\begin{figure}
  \inputminted{haskell}{usageO.hs}
  \caption{Usage example}
  \end{figure}
Clients of the library are somewhat more restricted, in the sense that they cannot create a
value of type \texttt{OrderedBy comp t} without going through the library's public API.

\begin{minted}{haskell}
minimum_O1 :: SortedBy comp [a] -> Maybe a
minimum_O1 xs = case (the xs) of
    []    -> Nothing
    (x:_) -> Just x
\end{minted}

\subsection{Conjuring a name}

Finally, for the user to be able to \emph{use} this library, there must be a way for
them to create \texttt{Named} values from normal values. The library must export a
function similar to this:

\begin{lstlisting}
name :: a -> (forall name. Named name a -> t) -> t
name x k = k (coerce x)
\end{lstlisting}

This function is quite similar to \texttt{sizing} from the previous section, and the rank-2
type gives it a bit of an ominous feel. You might wonder: why not just have a function
with a simple type like this?

\begin{lstlisting}
any_name :: a -> Named name a
any_name = coerce
\end{lstlisting}

The crux of the issue is all about \emph{who gets to choose} what \texttt{name} will be.
In the signature of \texttt{any\_name}, the \emph{caller} gets to select the types \texttt{a}
and \texttt{name}. In particular, they can attach any name they would like!

If that still does not sound so bad, consider this code:

\begin{minted}{haskell}
up, down :: Named () (Int -> Int -> Ordering)
up   = any_name (<)
down = any_name (>)

list1 = sortBy up   [1,2,3]
list2 = sortBy down [1,2,3]

merged = the (mergeBy up list1 list2) :: [Int]
-- [1,2,3,3,2,1]
\end{minted}

Now compare to the analogous program, using \texttt{name} instead of \texttt{any\_name}:
\begin{minted}{haskell}
name (<) $ \up ->
  name (>) $ \down ->
    let list1 = sortBy up   [1,2,3]
        list2 = sortBy down [1,2,3]
    in the (mergeBy up list1 list2)
\end{minted}

resulting in a compile-time error:

\begin{lstlisting}
  • Couldn't match type "name1" with "name"
        ...
    Expected type: SortedBy name [Integer]
      Actual type: SortedBy name1 [Integer]
\end{lstlisting}

A general rule of thumb for library authors is: \emph{a ghost should not appear in the return type,
  unless it also appears in an argument's type}. This simple rule ensures that
the user of the library will not be allowed to materialize ghosts out of thin air.

\section{Case Study \#2: \texttt{Maybe}-free lookup in containers}

\begin{filecontents*}{justified.hs}
newtype JMap φ k v = JMap (Map k v)
    deriving Functor

newtype JKey φ k = Element k

instance The (JMap φ k v) (Map k v)
instance The (JKey φ k)  k

member ::  k -> JMap φ k v -> Maybe (JKey φ k)

lookup   :: JKey φ k -> JMap φ k v -> v

reinsert
  :: JKey φ k -> v -> JMap φ k v -> JMap φ k v

withMap
  :: Map k v  -> (forall φ. JMap φ k v -> t) -> t
\end{filecontents*}

\begin{filecontents*}{justified-usage.hs}
test_table = Map.fromList [ (1, "Hello")
                          , (2, "world!") ]

withMap test_table $ \table ->
  case member 1 table of
 
    Nothing  -> putStrLn "Missing key!"
  
    Just key -> do
      putStrLn ("Found key: " ++ show (the key))
      putStrLn ("Value in map 1: " ++
                lookup key table)
      
      let table'  = reinsert key "Howdy" table
          table'' = fmap (map upper) table
      putStrLn ("Value in map 2: " ++
                lookup key table')
      putStrLn ("Value in map 3: " ++
                lookup key table'')
{- Output:
Found key: 1
Value in map 1: Hello
Value in map 2: Howdy
Value in map 3: HELLO
-}       
\end{filecontents*}

\begin{figure*}
  \begin{minipage}{0.48\textwidth}
    \inputminted{haskell}{justified.hs}
  \end{minipage}
  \begin{minipage}{0.48\textwidth}
    \inputminted{haskell}{justified-usage.hs}
  \end{minipage}
\end{figure*}

\subsection{Application: a type for directed graphs}

\begin{minted}{haskell}
data Neighbors phi = Neighbors
    { outEdges :: [JKey φ Vertex]
    , inEdges  :: [JKey φ Vertex] }
  
type Digraph φ = JMap φ Vertex (Neighbors φ)
\end{minted}

\begin{minted}{haskell}
addEdge   :: Vertex φ -> Vertex φ -> (forall φ'. Digraph φ' -> t) -> t
\end{minted}


\begin{minted}{haskell}
check :: Int -> Digraph φ -> Either (FreshVertex φ) (Vertex φ)
fresh :: Digraph φ -> FreshVertex φ
addVertex :: FreshVertex φ -> Digraph φ ->
            (forall φ'. (Vertex φ', Vertex φ -> Vertex φ', Digraph φ') -> t) -> t
\end{minted}

\subsection{Faster lookup}
Although \texttt{justified-containers} defines a simple \texttt{newtype} wrapper for
the key-plus-phantom-proof type, more interesting information about the location of
the key within the corresponding data structure can sometimes be attached.

For example, imagine a simple binary search tree backed by a vector of key-value pairs.
As in the previous example, we will give the \texttt{BST} type a phantom parameter that
represents the set of valid keys present in the tree. But instead of wrapping the key
type directly, we will use an index-plus-phantom-proof representation for keys.

\begin{lstlisting}[language=Haskell]
newtype BST phi k v = BST (Vector (k,v))
newtype Index phi   = Index Int

toBST :: Ord k => Vector (k,v) -> BST phi k v

find   :: Ord k => k -> BST phi k v -> Maybe (Index phi)
access :: Index phi -> BST phi k v -> (k,v)
\end{lstlisting}

\section{Case Study \#3: Encoding arbitrary properties}

\begin{figure*}
  \hspace{-2.5in}
  \scalebox{0.6}{{
  \begin{minipage}{\textwidth}
    \begin{prooftree}
      \AxiomC{}
      \RightLabel{\scriptsize (p)}\UnaryInfC{$\texttt{IsNil}(x) \wedge |x| = 0$}
      \UnaryInfC{$|x| = 0$}
      \UnaryInfC{$0 = |x|$}
      \AxiomC{}
      \RightLabel{\scriptsize (eq)}\UnaryInfC{$|x| = 1 + n$}
      \BinaryInfC{$0 = 1 + n$}
      \AxiomC{}
      \UnaryInfC{$\forall m \in \mathbb{N}.~ \neg (0 = 1 + m)$}
      \UnaryInfC{$\neg (0 = 1 + n)$}
      \BinaryInfC{$\bot$}
      \UnaryInfC{$\texttt{IsCons}(x)$}
      \RightLabel{\scriptsize (p)}\UnaryInfC{$\texttt{IsNil}(x) \wedge |x| = 0 \longrightarrow \texttt{IsCons}(x)$}
      \AxiomC{}
      \RightLabel{\scriptsize (q)}\UnaryInfC{$\texttt{IsCons}(x) \wedge |x| = 1 + |\texttt{Tail}(x)|$}
      \UnaryInfC{$\texttt{IsCons}(x)$}
      \RightLabel{\scriptsize (q)}\UnaryInfC{$\texttt{IsCons}(x) \wedge |x| = 1 + |\texttt{Tail}(x)| \longrightarrow \texttt{IsCons}(x)$}
      \AxiomC{}
      \UnaryInfC{$\forall \ell \in \texttt{List}_T.~ (\texttt{IsNil}(\ell) \wedge |\ell| = 0) \vee (\texttt{IsCons}(\ell) \wedge |\ell| = 1 + |\texttt{Tail}(\ell)|)$}
      \UnaryInfC{$(\texttt{IsNil}(x) \wedge |x| = 0) \vee (\texttt{IsCons}(x) \wedge |x| = 1 + |\texttt{Tail}(x)|)$}
      \TrinaryInfC{$\texttt{IsCons}(x)$}
      \RightLabel{\scriptsize (eq)}\UnaryInfC{$|x| = 1 + n \longrightarrow \texttt{IsCons(x)}$}
    \end{prooftree}
  \end{minipage}
    }
  }
  \caption{A proof that lists with nonzero length satisfy the \texttt{IsCons} predicate,
      in natural deduction style. Compare with the same proof using the \texttt{Proof} monad
      in listing ****; the steps after the \texttt{(|/)} operator correspond to the leftmost
      deductions in this proof tree. Note a slight difference: the listing proves
      $|x| = 1 + n \vdash \texttt{IsCons}(x)$, while the derivation in this figure
      proves $\vdash |x| = 1 + n \longrightarrow \texttt{IsCons}(x)$.}
\end{figure*}

\begin{filecontents*}{nzlic.hs}
nonzero_length_implies_cons
  :: (Length xs == Succ n)
  -> Proof (IsCons xs)

nonzero_length_implies_cons eq =
  do  toSpec length
   |$ or_elimR and_elimL
   |/ and_elimR
   |. symmetric
   |. transitive' eq
   |. (contradicts' $$ zero_not_succ)
   |. absurd
\end{filecontents*}

\inputminted{haskell}{nzlic.hs}

%% \section{Case Study \#2: A stack-based calculator}

%% In this case study, we investigate how the GDP technique can be used to
%% acheive some of the safety and expressivity features of dependent types,
%% within a system that only supports non-dependent types.

%% To showcase the method, we will develop a library for carrying out computations
%% with a stack-based calculator. The basic state of the calculator will be
%% represented as a stack of \texttt{Double}s, but we will also introduce a
%% phantom parameter: the ghost of the stack's size.

%% \begin{lstlisting}[language=Haskell]
%%   newtype Stack $nu$ = Stack [Double]
%% \end{lstlisting}

%% We also introduce a type constructor \texttt{S} with the intention that
%% if $\nu$ represents a stack of height $n$, then $\texttt{S } \nu$ represents
%% a stack of height $n + 1$. Since \texttt{S} is only supposed to be applied
%% to phantom types, there is no need for any type of the form \texttt{S t} to
%% be inhabited. As a result, we can use an empty data declaration:
%% \begin{lstlisting}[language=Haskell]
%%   data S $nu$    -- empty data declaration, an uninhabitable type
%% \end{lstlisting}
%% The safe stack operations can then be encoded:
%% \begin{lstlisting}[language=Haskell]
%%   pushStack :: Double -> Stack $nu$ -> Stack (S $nu$)
%%   pushStack x (Stack stk) = Stack (x:xs)

%%   popStack :: Stack (S $nu$) -> (Double, Stack $nu$)
%%   popStack (Stack (x:xs)) = (x, Stack xs)
%%   popStack (Stack []) = error "the spirits have been violated"
%% \end{lstlisting}
%% Finally, we should provide a method for generating ghosts:
%% \begin{lstlisting}[language=Haskell]
%%   testStack :: Stack $nu$ -> (forall $mu$. Maybe (Stack (S $mu$)) -> t) -> t
%%   testStack stack action = action $ case stack of
%%       Stack (_:_) -> Just (coerce stack)
%%       _           -> Nothing
%% \end{lstlisting}

%% \subsection{Impredicativity}
%% There is an interesting issue related to impredacative types that can be seen
%% in the encoding of \texttt{testStack}. First, \texttt{testStack} should take a
%% general stack as input, and either reify the fact that there is at least one
%% element, or result in \texttt{Nothing}. So the natural type to expect is
%% \[\texttt{Stack } \nu \to \texttt{Maybe } (\exists \mu.~ \texttt{Stack } (\texttt{S } \mu)).\]
%% But a straightforward encoding as a rank-2 type with universal quantification leads to
%% \[\texttt{Stack } \nu \to \texttt{Maybe } ((\forall \mu.~ \texttt{Stack } (\texttt{S } \mu) \to t) \to t).\]
%% *****
%% Instead, we have to use the weaker encoding
%% \[\texttt{Stack } \nu \to \exists \mu.~ \texttt{Maybe } (\texttt{Stack } (\texttt{S } \mu)),\]
%% corresponding to the rank-2 type
%% \[\texttt{Stack } \nu \to (\forall \mu.~ \texttt{Maybe } (\texttt{Stack } (\texttt{S } \mu)) \to t) \to t.\]

%% \begin{lstlisting}[language=Haskell]
%%   iget :: IxMonadStateT m p p p
%%   iput :: q -> IxMonadStateT m p q ()
%%   imodify :: (p -> q) -> IxMonadStateT m p q ()

%%   return :: a -> IxMonadStateT m p p a
%%   (>>=) :: IxMonadStateT m p q a -> (a -> IxMonadStateT m q r b) -> IxMonadStateT m p r b
%% \end{lstlisting}

%% \begin{lstlisting}[language=Haskell]
%%   type StackOp $nu$ $mu$ t = IxMonadStateT IO (Stack $nu$) (Stack $mu$) t

%%   push :: Double -> StackOp $nu$ (S $nu$) ()
%%   push x = imodify (pushStack x)
  
%%   pop :: StackOp (S $nu$) $nu$ Double
%%   pop = do
%%     stack <- iget
%%     let (top, rest) = popStack stack
%%     iput rest
%%     return top
%% \end{lstlisting}

%% \section{Case Study \#3: Liquid Haskell inside}

%% \section{Case Study \#4: Ghosts of dependent types}

%% \begin{lstlisting}[language=Haskell]
%%   data S $nu$
%% \end{lstlisting}

%% \begin{lstlisting}[language=Haskell]
%%   pf_commutes :: Plus n m -> Plus m n
%%   pf_assocL   :: Plus a (Plus b c) -> Plus (Plus a b) c
%%   pf_assocR   :: Plus (Plus a b) c -> Plus a (Plus b c)
%%   pf_ident    :: Plus n Z -> n
%%   pf_destruct :: n -> Either (n :=: Z) (n :=: S m)
%% \end{lstlisting}

%% \begin{lstlisting}[language=Haskell]
%%   push :: a -> Stack n a -> Stack (S n) a
%%   pop  :: Stack (S n) a -> (a, Stack n a)
%%   test :: Stack n a -> Either (Stack Z) (Stack (S m))
%% \end{lstlisting}

%% \begin{lstlisting}[language=Haskell]
%%   opAdd :: Stack (S (S n)) -> Stack (S n)
%%   opAdd stack = push (x + y) stack''
%%     where
%%       (x, stack')  = pop stack
%%       (y, stack'') = pop stack'
%% \end{lstlisting}


%% \begin{lstlisting}[language=Haskell]
%%   data Column rho
%%   data Row kappa
%%   data Matrix rho kappa
%% \end{lstlisting}

%% \begin{lstlisting}[language=Haskell]
%% class Euclidean v where
%%   (.+.) :: v -> v -> v
%%   (**)  :: Double -> v -> v
%%   zero  :: v
%%   type family Transposed :: *
%%   transpose :: v -> vT


%% instance Euclidean (Col kappa) where
%%   (Col v1) .+. (Col v2) = Col (v1 Vector.(.+.) v2)
%%   zero = Col (Vector.zero)
%%   type Transposed (Col kappa) = Row kappa
%%   transpose = coerce

%% instance Euclidean (Row rho) where
%%   (Row v1) .+. (Row v2) = Row (v1 Vector.(.+.) v2)
%%   zero = Row (Vector.zero)
%%   type Transposed (Row rho) = Col rho
%%   transpose = coerce

%% instance Euclidean (Matrix rho kappa) where
%%   type Transposed (Matrix rho kappa) = Matrix kappa rho
%%   transpose (Matrix m) = Matrix (V.tranpose m)
  
%% transpose :: Matrix rho kappa -> Matrix kappa rho
%% (.+.) :: Matrix rho kappa -> Matrix rho kappa -> Matrix rho kappa
%% (.*.) :: Matrix rho kappa -> Matrix kappa gamma -> Matrix rho gamma
%% (*.) :: Matrix rho kappa -> Column rho -> Column kappa
%% (.*) :: Row kappa -> Matrix rho kappa -> Row rho
%% inner :: Row kappa -> Col kappa -> Double
%% outer :: Col kappa -> Row rho   -> Matrix rho kappa

%% \end{lstlisting}

%% \begin{lstlisting}[language=Haskell]
%% dimRow :: Vector Double -> (forall rho. Row rho -> t) -> t
%% dimCol :: Vector Double -> (forall kappa. Col kappa -> t) -> t
%% dimMat :: Matrix Double -> (forall rho kappa. Mat rho kappa -> t) -> t

%% alignRow :: Row rho -> Row rho' -> Maybe (rho <~> rho')
%% \end{lstlisting}

%% \begin{lstlisting}[language=Haskell]
%% phi :subsetX: phi' = ProperSubset
%%   { pf_subset   :: Vertex phi -> Vertex phi'
%%   , pf_newvtx   :: Vertex phi'
%%   , pf_classify :: Vertex phi' -> Maybe (Vertex phi)
%%   }
%% \end{lstlisting}

%% \begin{lstlisting}[language=Haskell]
%% coerceWith :: (a ~> b) -> (a -> b)

%% coerceToR :: (a <~> b) -> (a -> b)
%% coerceToL :: (a <~> b) -> (b -> a)

%% phi :subsetX: phi' = ProperSubset
%%   { pf_subset   :: Vertex phi ~> Vertex phi'
%%   , pf_newvtx   :: Vertex phi'
%%   , pf_classify :: Vertex phi' -> Maybe (Vertex phi)
%%   }
%% \end{lstlisting}


\end{document}
